\qquad A partir da implementação e dos experimentos realizados com os
cinco algoritmos, foi possível observar, na prática, diferenças de
comportamento que confirmam a análise teórica de cada um. Os tempos
citados a seguir referem-se a instâncias de 1.000.000 de elementos.

\qquad O \textbf{Insertion Sort} e o \textbf{Bubble Sort} compartilham uma
característica importante: ambos são sensíveis ao grau de ordenação
da entrada, apresentando desempenho linear ($O(n)$) quando os dados já
estão ordenados (0,0006 s e 0,0003 s, respectivamente) e desempenho
quadrático ($O(n^2)$) nos casos médio e pior (entradas randômica e
decrescente). Entre os dois, o resultado depende do tipo de entrada:
com a entrada decrescente os tempos foram equivalentes (167,0 s para o
Insertion Sort e 155,6 s para o Bubble Sort), mas com a entrada
randômica o Bubble Sort foi cerca de 11,6 vezes mais lento (938,7 s
contra 80,6 s). Parte dessa diferença decorre do dobro de comparações do
Bubble Sort em uma entrada aleatória ($n(n-1)/2$ contra $n(n-1)/4$) e,
provavelmente, de efeitos de hardware que a contagem de operações não
captura.

\qquad O \textbf{Selection Sort} se destacou por um comportamento distinto:
sua complexidade é $\Theta(n^2)$ em \textit{todos} os casos (melhor, médio e
pior), já que o número de comparações não depende da ordem dos dados
--- só o de trocas varia, mas é sempre $O(n)$, irrelevante diante das
$O(n^2)$ comparações.
Isso foi confirmado experimentalmente: os tempos medidos para as
entradas crescente, decrescente e randômica ficaram praticamente
idênticos entre si (932,5 s a 937,3 s, uma variação de 0,5\%). Por não
se beneficiar de entradas ordenadas, foi o mais lento na entrada
crescente (934,9 s, contra 0,0006 s do Insertion Sort) e equivalente ao
Bubble Sort na randômica. Sua vantagem
é o número previsível e mínimo de trocas (no máximo $n-1$), útil quando
a operação de escrita é cara.

\qquad O \textbf{Shell Sort} apresentou desempenho muito superior ao dos
três algoritmos anteriores em todos os tipos de entrada testados
(0,008 s, 0,012 s e 0,106 s para as entradas crescente, decrescente e
randômica), mesmo utilizando a sequência de intervalos mais simples (a
original de Shell). Embora essa sequência tenha pior caso teórico
$\Theta(n^2)$, ele exige uma disposição de dados construída
especificamente para isso, o que não ocorre nem com uma entrada
totalmente decrescente --- por isso as entradas crescente e decrescente
exibiram comportamento próximo de $O(n \log n)$, e a randômica (sem
fórmula fechada simples de caso médio) foi cerca de 760 vezes mais
rápida que o Insertion Sort.

\qquad O \textbf{Merge Sort} tem complexidade $\Theta(n \log n)$ em todos os
casos, independentemente da ordem inicial dos dados, e foi o único
algoritmo cujo tempo não apresentou grande variação com o tipo de
entrada: 0,021 s (crescente), 0,021 s (decrescente) e 0,078 s
(randômica). Para a entrada randômica foi o mais rápido dos cinco
(0,078 s, contra 0,106 s do Shell Sort e 80,6 s do Insertion Sort, o
melhor dos três algoritmos quadráticos nesse cenário, ou seja, cerca de
mil vezes mais rápido). Quando a entrada já é crescente ou decrescente,
o Shell Sort foi mais rápido que o Merge Sort (0,008 s e 0,012 s contra
0,021 s), e o Insertion e o Bubble Sort, lineares na entrada crescente,
ainda mais. O preço do Merge Sort é a memória
adicional proporcional ao tamanho da entrada (um vetor auxiliar de
cerca de 4 MB para 1.000.000 de inteiros), enquanto os outros quatro
algoritmos ordenam no próprio vetor.

\qquad Em resumo, os experimentos evidenciam que a escolha do algoritmo
importa mais do que qualquer ajuste de detalhe: para 1.000.000 de
elementos randômicos, os algoritmos de menor complexidade (Merge e
Shell Sort) levaram cerca de 0,1 s, contra 80 a 940 s dos três
algoritmos $O(n^2)$ --- uma diferença de três a quatro ordens de
grandeza. Entre os algoritmos $O(n^2)$, o Insertion Sort foi o mais
regular: nunca foi o pior em nenhum cenário e foi, de longe, o melhor
na entrada randômica, sendo preferível quando se espera entradas quase
ordenadas; o Selection Sort tem a vantagem de um número mínimo de
trocas, e o Bubble Sort, apesar de didático, não apresentou vantagem
prática clara sobre os demais, a não ser na entrada crescente, onde
foi o mais rápido por uma margem muito pequena. O Merge Sort é a escolha mais
segura quando se precisa de desempenho garantido para qualquer
entrada e há memória disponível; o Shell Sort, que mostra como uma
modificação simples do Insertion Sort (intervalos decrescentes) produz
um ganho expressivo, é uma boa alternativa quando a memória é restrita
ou os dados tendem a estar parcialmente ordenados.

\qquad Cada instância foi executada uma única vez, com instâncias
aleatórias distintas por algoritmo, e medida com \texttt{clock()}, de
resolução limitada; por isso, diferenças pequenas, como a de 7\% entre
o Insertion e o Bubble Sort na entrada decrescente, não indicam
superioridade de um algoritmo, ao contrário das diferenças de ordens
de grandeza, que são robustas.
